\documentclass[10pt,a4paper]{amsart}
\usepackage[margin=19mm]{geometry}
\usepackage{amsmath,amssymb,mathtools}
\usepackage[scr=rsfs,cal=euler]{mathalfa}
\usepackage{xfrac}
\usepackage[unicode,hidelinks,bookmarks=false]{hyperref}
\newcommand{\ie}{i.e.\ }
\newcommand{\cf}{cf.\ }
\newcommand{\resp}{respectively\ }
\newcommand{\Subsection}{\subsection}
\providecommand{\nobreakdash}{\nobreak\hbox{-}}
\setlength{\emergencystretch}{2em}
\multlinegap0pt
\usepackage{bm}
\usepackage{amssymb}
\usepackage{xcolor}
\newtheorem{proposition}{Proposition}
\def\H{\bm{H}}
\def\L{\bm{L}}
\def\vv{\bm{v}}
\title{\emph{On a Thermodynamically Consistent Diffuse-Interface Model for Incompressible Two-Phase Flows \\ with Chemotaxis and Mass Transport }}
\author{Andrea Giorgini$^\dagger$, \ \ \ Jingning He$^\ddagger$, \ \ \ Hao Wu$^\ast$}
\date{Source: arXiv:2512.22585v1 (CC BY 4.0)}
\begin{document}
\maketitle

\noindent\emph{Source and licence.} \emph{On a Thermodynamically Consistent Diffuse-Interface Model for Incompressible Two-Phase Flows \\ with Chemotaxis and Mass Transport }. Version arXiv:2512.22585v1 (CC BY 4.0), distributed by arXiv under the Creative Commons Attribution 4.0 International licence, http://creativecommons.org/licenses/by/4.0/, as recorded in the arXiv licence metadata for this exact version and shown on the abstract page https://arxiv.org/abs/2512.22585v1. Full licence notice: \texttt{licenses/arxiv-2512.22585-CC-BY-4.0.txt}.\par\smallskip

\noindent\textit{Editorial scope.} Counted unit: the article's proposition sigma-strong (source0.tex lines 3488--3510). The article's complete original proof of it is delivered with it (source0.tex lines 3512--3579, one \texttt{proof} environment of 68 lines). No mathematics is written, repaired or re-derived: the statement and the proof are the article's own lines, in the article's own notation, and no retained line is substituted or re-flowed.  Retained uncounted as the source-local context the proof needs: lines 3375--3391; lines 3427--3430.  Nothing was omitted from any retained span, so the package carries no omission record.  No retained line contains a citation, so the wrapper carries no bibliography and no bibliography entry.  No retained line contains a figure environment, an image-inclusion command or drawing code, so no image is delivered and the package declares no asset dependency.  Editorial formatting only, in addition: an AMS \texttt{amsart} preamble that restates the article's own declarations and macros used by the retained lines, namely the environments \texttt{proposition} on separate counters, together with the standard packages those lines need.  The retained environments print in this document as Proposition 1 in order of appearance, whereas the article prints the same items with the numbering of its own full document.  The complete native source of the article as distributed in the arXiv e-print for this exact version is bundled unchanged as \texttt{sources/191746/source0.tex}.
Let $\vv$, $\varphi$ be two given functions with suitable regularity properties. Consider the following diffusion equation with convection:
%
\begin{alignat}{3}
&\partial_t\sigma +\vv\cdot\nabla\sigma -\Delta\sigma
=  \mathrm{div}\big(\beta'(\varphi)\sigma \nabla \varphi\big),
\quad \text{in}\ \Omega\times (0,\infty),
\label{re-di-sigma}
\end{alignat}
%
equipped with the following boundary and initial conditions
%
\begin{alignat}{3}
&\partial_{\boldsymbol{n}} \sigma =0  \quad \text{on}\ \partial\Omega\times(0,\infty),
\qquad \sigma|_{t=0}=\sigma_{0} \quad \text{in}\ \Omega.
\label{re-di-sigma-bdini}
\end{alignat}
%

\begin{align}
\int_\Omega\sigma(t)\,\mathrm{d}x=\int_\Omega \sigma_0\,\mathrm{d}x,\quad \forall\, t\geq 0.
\label{conv-sig}
\end{align}

\begin{proposition}
\label{sigma-strong}
Let $\Omega$ be a bounded domain in $\mathbb{R}^2$ with a $C^2$ boundary. Assume that (H4) is satisfied and
%
\begin{align*}
&\vv \in L^{\infty}(0,\infty;\L^2_{\sigma}(\Omega))\cap L^2(0,\infty;\H^1_{0,\sigma}(\Omega)),
\\
&\varphi \in L^\infty(0,\infty;H^1(\Omega))\cap L^2_{\mathrm{uloc}}([0,\infty);W^{2,q}(\Omega)\cap H^2_N(\Omega)),
\end{align*}
%
for some $q>2$. Then for any initial datum $\sigma_0 \in H^1(\Omega)$ satisfying $\sigma_0 \geq 0$ almost everywhere in $\Omega$, problem
\eqref{re-di-sigma}--\eqref{re-di-sigma-bdini} admits a unique global strong solution $\sigma$ in $\Omega\times [0,\infty)$ such that
%
\begin{equation}
\label{REG1}
\begin{split}
\sigma \in L^{\infty}(0,T;H^1(\Omega)) \cap L^2(0,T; H^2(\Omega))\cap H^1(0,T;L^2(\Omega)),
\end{split}
\end{equation}
%
for any $T>0$, with $\sigma(x,t)\geq 0$ almost everywhere in $\Omega \times (0,\infty)$,
which satisfies \eqref{re-di-sigma} almost everywhere in $\Omega\times(0,\infty)$. Moreover, we have $\partial_{\boldsymbol{n}} \sigma =0$ almost everywhere on $\partial\Omega\times(0,\infty)$ and  $\sigma|_{t=0}=\sigma_{0}$ almost everywhere in $\Omega$.
\end{proposition}

\begin{proof}
We only derive the necessary higher-order \textit{a priori} estimates.
Multiplying \eqref{re-di-sigma} by $-\Delta\sigma$ and integrating over $\Omega$, we obtain
%
\begin{equation}
\begin{aligned}
\frac{1}{2}\frac{\mathrm{d}}{\mathrm{d} t}\|\nabla \sigma\|^2+\left\|\Delta\sigma\right\|^2
&=\int_{\Omega}\vv \cdot \nabla \sigma\Delta\sigma\,\mathrm{d}x
- \int_{\Omega}\mathrm{div}\big(\beta'(\varphi)\sigma \nabla \varphi\big) \Delta\sigma \,\mathrm{d}x.
\end{aligned}
\notag
\end{equation}
%
Using H\"older's inequality, Young's inequality, the Sobolev embedding theorem, the Gagliardo--Nirenberg inequality, the Poincar\'e--Wirtinger inequality and \eqref{conv-sig}, we have
%
\begin{equation*}
\begin{aligned}
\left|\int_{\Omega}\vv \cdot \nabla \sigma \Delta\sigma \,\mathrm{d}x\right|
&\le \|\vv \|_{\L^4(\Omega)}\|\nabla \sigma\|_{\bm{L}^4(\Omega)}\|\Delta\sigma\|\\
&\le \frac18\|\Delta\sigma\|^2 + C\|\vv \|\|\vv \|_{\bm{H}^1(\Omega)}
\|\nabla \sigma\| (\|\Delta\sigma\| + \overline{\sigma_0}) \\
&\le \frac14\|\Delta\sigma\|^2 + C\|\nabla \vv \|^2\|\nabla \sigma\|^2+C,
\end{aligned}
\end{equation*}
and
\begin{equation*}
\begin{aligned}
&\left|\int_{\Omega}\mathrm{div}\big(\beta'(\varphi)\sigma \nabla \varphi\big) \Delta\sigma \,\mathrm{d}x\right|
\\
&\quad\le \|\beta'(\varphi)\|_{L^\infty(\Omega)}\|\nabla \sigma\|\|\nabla \varphi\|_{\bm{L}^\infty(\Omega)}\|\Delta\sigma\|
+ \|\beta''(\varphi)\|_{L^\infty(\Omega)}\|\sigma\|_{L^\frac{2q}{q-2}(\Omega)} \|\nabla \varphi\|^2_{\bm{L}^{2q}(\Omega)}\|\Delta\sigma\|
\\
&\qquad +\|\beta'(\varphi)\|_{L^\infty(\Omega)}\|\sigma\|_{L^\frac{2q}{q-2}(\Omega)} \|\Delta \varphi\|_{L^q(\Omega)}\|\Delta\sigma\| \\
&\quad \leq C \|\nabla \sigma\|\|\varphi\|_{W^{2,q}(\Omega)}\|\Delta \sigma\|
+ C
\| \sigma \|_{H^1(\Omega)}
(\| \nabla \varphi\| +1)
 \| \varphi\|_{W^{2,q}(\Omega)} \| \Delta \sigma\|
\\
&\quad\le \frac14\|\Delta\sigma\|^2
+ C\|\varphi\|_{W^{2,q}(\Omega)}^2(\|\nabla \sigma\|^2+1).
\end{aligned}
\end{equation*}
%
From the above estimates, we can derive the following differential inequality
%
\begin{equation}
\begin{aligned}
\frac{\mathrm{d}}{\mathrm{d} t} \|\nabla \sigma\|^2 + \|\Delta\sigma\|^2
&\le C\big(\|\nabla \vv \|^2+\|\varphi\|_{W^{2,q}(\Omega)}^2\big) \|\nabla \sigma\|^2
+ C\big(\|\varphi\|_{W^{2,q}(\Omega)}^2+1\big).
\end{aligned}
\label{es-sigH0}
\end{equation}
%
Applying Gronwall's lemma and the Poincar\'{e}--Wirtinger inequality, we obtain
%
\begin{align}
\|\sigma\|_{L^\infty(0,T;H^1(\Omega))}\leq K_3,\quad \|\sigma\|_{L^2(0,T;H^2(\Omega))}\leq K_3,
\label{es-sigH1}
\end{align}
%
for any $T>0$, where the constant $K_3>0$ depends on $\|\sigma_0\|_{H^1(\Omega)}$,
$\|\vv\|_{L^\infty(0,\infty;\L^2_{\sigma}(\Omega))}$, $\|\vv\|_{L^2(0,\infty;\H^1_{0,\sigma}(\Omega))}$,
$\|\varphi \|_{L^\infty(0,T;H^1(\Omega))}$, $\|\varphi\|_{L^2(0,T;W^{2,q}(\Omega))}$, and $T$. Finally, by a comparison in the equation \eqref{re-di-sigma}, it is straightforward to check that $\partial_t\sigma \in L^2(0,T;L^2(\Omega))$.

The proof of Proposition \ref{sigma-strong} is complete.
\end{proof}

\end{document}
